\chapter{Discussão e Conclusão}

\section{Nível de Complexidade do Problema}

Intuitivamente, a complexidade do problema decorre da necessidade de decidir simultaneamente a alocação e o roteamento de fluxos em uma rede com múltiplos níveis.

No âmbito da Teoria da Complexidade Computacional, o problema é classificado como NP-difícil, uma vez que se trata de um Problema de Programação Linear Inteira e para problemas desta natureza não se conhece nenhum algoritmo capaz de resolvê-los em tempo polinomial no pior caso.

No pior caso, o algoritmo pode ser forçado a explorar praticamente todas as combinações possíveis de soluções. Como o número dessas combinações cresce exponencialmente com o tamanho da instância, o tempo de execução pode apresentar comportamento exponencial, isto é:

\[
\text{tempo} \sim O(2^n) \quad \text{(ou pior)}
\]

onde $n$ representa o tamanho da instância. No contexto da Programação Linear Inteira, esse tamanho está em função do limite do espaço de busca explorado pelos algoritmos de otimização (como o \textit{Branch-and-Bound}) e esse limite é ditado pelo número total de variáveis de decisão do sistema.

O Modelo 2 (sem penalidade) foi idealizado sob uma ótica mais realista que os demais modelos e, por isso, é válido analisar a complexidade computacional gerada por essa abordagem.

O Modelo 2 exige a operação com três tipos de variáveis de decisão tridimensionais ($x_{ikv}$, $y_{kjv}$ e $z_{jpv}$) e uma variável bidimensional de falta ($f_{pv}$). Como consequência, a quantidade de vezes que as variáveis aparecem na função objetivo e o tamanho da matriz de restrições crescem de forma multiplicativa em relação ao tamanho dos conjuntos de municípios de origem ($|I|$), municípios de destino ($|J|$), centros estaduais ($|K|$), postos de vacinação ($|P(j)|$) e tipos de vacina ($|V|$).

Isso torna este modelo impossível de ser resolvido manualmente e altamente exigente do ponto de vista computacional. Mensurando a complexidade gerada por essa modelagem, tem-se:

\textbf{Quantidade de vezes que as variáveis de decisão aparecem na função objetivo:}
\[
(|I| \cdot |K| \cdot |V|)
+
(|K| \cdot |J| \cdot |V|)
+
(|P| \cdot |V|)
+
(|P| \cdot |V|)
\text{ vezes}
\]

\textbf{Tamanho da matriz de restrições (Soma de todas as regras do sistema):}
\begin{align*}
& (|I| \cdot |V|) \quad && \text{(Oferta Intermunicipal)} \\
+ \;& (|K| \cdot |V|) \quad && \text{(Conservação do Fluxo no Estado)} \\
+ \;& (|P| \cdot |V|) \quad && \text{(Demanda nos Postos de Saúde)} \\
+ \;& (|I| \cdot |K|) + (|K| \cdot |J|) + (|P|) \quad && \text{(Capacidades Físicas de Transporte)} \\
+ \;& (|K| \cdot |J| \cdot |V|) \quad && \text{(Consistência Estadual)} \\
+ \;& (|J| \cdot |V|) \quad && \text{(Balanço de Fluxo Municipal)} \\
+ \;& (|I| \cdot |K| \cdot |V|) + (|P| \cdot |V|) \quad && \text{(Viabilidade Temporal)}
\end{align*}

Para o estudo de caso do cenário nacional utilizado neste trabalho, foram mapeados 27 municípios de origem ($|I|=27$), 27 municípios de destino ($|J|=27$), 27 centros estaduais ($|K|=27$), 54 postos de vacinação globais ($|P|=54$, sendo 2 por capital) e 4 tipos de vacina ($|V|=4$).

Aplicando essas dimensões práticas à estrutura das equações, foram definidos os seguintes valores exatos para o modelo:

\textbf{Quantidade de vezes que as variáveis de decisão aparecem na função objetivo:}
\[
(27 \cdot 27 \cdot 4) 
+ 
(27 \cdot 27 \cdot 4) 
+ 
(54 \cdot 4) 
+
(54 \cdot 4)
\]

\[
= 2916 + 2916 + 216 + 216 = 6.264 \text{ vezes}
\]

\textbf{Quantidade de linhas na matriz de restrições:}
\begin{align*}
& (27 \cdot 4) \quad && \text{(Oferta Intermunicipal)} \\
+ \;& (27 \cdot 4) \quad && \text{(Conservação do Fluxo no Estado)} \\
+ \;& (54 \cdot 4) \quad && \text{(Demanda nos Postos de Saúde)} \\
+ \;& (27 \cdot 27) + (27 \cdot 27) + (54) \quad && \text{(Capacidades Físicas de Transporte)} \\
+ \;& (27 \cdot 27 \cdot 4) \quad && \text{(Consistência Estadual)} \\
+ \;& (27 \cdot 4) \quad && \text{(Balanço de Fluxo Municipal)} \\
+ \;& (27 \cdot 27 \cdot 4) + (54 \cdot 4) \quad && \text{(Viabilidade Temporal)}
\end{align*}

\[
= 108 + 108 + 216 + 1512 + 2916 + 108 + 3132 = 8.100 \text{ linhas}
\]


\textbf{Tamanho da Instância:}

O tamanho da instância é a soma total das variáveis de decisão inteiras geradas pelos dados de entrada, ou seja, 6.264.

Logo, no pior caso, o tempo de execução pode ser:

\[
\text{tempo} \sim O(2^{6264}) \quad \text{(ou pior)}
\]

\textbf{Tamanho total da matriz do problema (Variáveis $\times$ Restrições):}
\[
6.264 \times 8.100 = 50.738.400 \text{ células}
\]


Esse número torna colossal o nível de complexidade gerado por essa modelagem, o que afeta diretamente o planejamento logístico:
\begin{itemize}
    \item \textbf{Dependência do Auxílio Computacional:} É fisicamente e humanamente impossível delinear um plano eficiente de redistribuição de vacinas sem o auxílio do poder computacional;
    \item \textbf{Necessidade de Ferramentas de Ponta:} O tamanho da matriz justifica o uso de \textit{solvers} avançados, como o Gurobi Optimizer, que garante a otimalidade global dentro de um tempo hábil.
\end{itemize}

\section{Limitações da Modelagem}

Para garantir a tratabilidade computacional foram adotadas simplificações que distanciam a simulação de um cenário logístico real em sua totalidade. Dessa forma, a modelagem possui uma série de limitações, como:

\begin{itemize}
    \item \textbf{Negligência de Instâncias Administrativas Intermediárias:} O modelo adota uma estrutura de rede condensada composta por, no máximo, três níveis principais de fluxo de suprimentos ($I \to K \to J \to P$). Contudo, na estrutura organizacional real do Sistema Único de Saúde (SUS), estados de grande extensão territorial, como o de Minas Gerais, podem operar com um nível adicional composto pelas Superintendências Regionais de Saúde (SRS), configurando um fluxo do tipo $I \to K \to \text{SRS} \to J \to P$;
    \item \textbf{Homogeneidade da Malha Rodoviária:} Na vida real, um eixo logístico entre São Paulo e Rio de Janeiro possui uma malha rodoviária infinitamente mais robusta do que a de um trajeto ligando municípios da Região Norte, que pode apresentar trechos rodoviários precários. Essa heterogeneidade da malha rodoviária que varia entre duas localidades não foi contemplada no modelo;
    \item \textbf{Simetria da Fórmula de Haversine:} A utilização da fórmula de Haversine acrescida de um fator de tortuosidade ($\beta$) cria um grafo simétrico, onde o custo de ida é exatamente igual ao custo de volta. Em redes rodoviárias reais, o relevo (como a subida ou descida de uma serra), a existência de vias de mão única, pedágios e rotas alternativas fazem com que o tempo de deslocamento seja assimétrico;
    \item \textbf{Incapacidade de Previsão do Trânsito Rodoviário:} O modelo considera a velocidade média do modal rodoviário ($\theta_{rod}$) como um parâmetro constante. Porém, na realidade urbana, o trânsito é um sistema dinâmico e imprevisível, pois está suscetível a intempéries climáticas, acidentes, bloqueios em rodovias, obras e grandes eventos (como partidas de futebol ou festividades), que podem travar o trânsito, impactando drasticamente o tempo de deslocamento $\tau_{ab}^{rod}$;
    \item \textbf{Agregação das Unidades de Vacina do Tipo V:} Um município real não possui um único prazo de validade para todas as suas vacinas de um tipo $v$, mas sim múltiplas unidades de vacina com datas de expiração distintas. O modelo não rastreia unidades individuais, mas sim um lote agregado de vacinas do tipo $v$ com um prazo de validade médio estimado ($e_{iv}$ ou $e_{jv}$). Essa agregação simplifica a modelagem, mas perde a precisão do rastreamento de unidades individuais;
    \item \textbf{Simplificação do Tempo do Modal Aéreo:} Na prática, o uso do modal aéreo compreende, nesta ordem, os tempos de transporte rodoviário do município de origem até o aeroporto, o tempo de voo entre aeroportos e o transporte rodoviário do aeroporto de destino até o município final. Porém, na aplicação e na modelagem, o tempo do modal aéreo ($\tau_{ab}^{aereo}$) considera somente o tempo de voo entre municípios.
\end{itemize}

\section{Limitações da Aplicação dos Modelos}

Embora o modelo matemático desenvolvido apresente solidez teórica para a otimização logística, a sua aplicação no cenário real do Sistema Único de Saúde (SUS) esbarra em desafios práticos, como:

\begin{itemize}
    \item \textbf{Sigilo dos Dados:} Informações sobre estoques e demandas de vacinas e datas de validade dos lotes não são dados públicos e são tratados como dados sensíveis pelas secretarias de saúde;
    \item \textbf{Atualização dos Dados:} A aplicação do modelo não depende somente da obtenção dos dados sigilosos, mas também que eles estejam 100\% atualizados, o que é um desafio dado o número massivo de lotes de vacinas por todo o país;
    \item \textbf{Volume Massivo de Dados:} O Brasil possui mais de 5.000 municípios e dezenas de milhares de postos de saúde, o que torna a coleta e a organização manuais desses dados uma tarefa humanamente impraticável. Por isso, se faz necessária uma coleta e organização automatizada desses dados.
\end{itemize}

Essas barreiras práticas justificam e reforçam a geração sintética dos dados e as simplificações adotadas neste estudo de caso. 

Ao reduzir a abstração computacional para as 27 capitais brasileiras e limitar a amostragem de postos de atendimento, foi possível comprovar o funcionamento lógico do algoritmo e analisar seus resultados. 

Conclui-se que o verdadeiro desafio da distribuição ótima de vacinas não reside apenas no processamento computacional das rotas ótimas, mas também na engenharia de dados necessária para coletar e organizar dados atualizados e em quantidade massiva.

\section{Relevância do Trabalho e Conclusões Finais}

Primeiramente, ressalta-se que toda a estruturação da modelagem matemática, bem como a estratégia de extração e tratamento dos dados, foram concebidas e desenvolvidas integralmente pelo autor. Este trabalho não consiste em uma derivação ou mera modificação de artigos presentes na bibliografia consultada, consolidando-se como uma contribuição autoral.

A pesquisa cumpriu com êxito o seu objetivo central: modelar, implementar e resolver o problema de redistribuição de vacinas em um cenário pandêmico no Brasil sob a ótica da Programação Linear Inteira, buscando minimizar o tempo logístico de transporte, ampliar o atendimento da demanda vacinal e contribuir para respostas mais eficientes em contextos de emergência sanitária.

Por meio da modelagem matemática, níveis de cadeia logística e a burocracia foram traduzidos com precisão por meio de variáveis de decisão e restrições, como a restrição da conservação do fluxo e a de consistência estadual. Essa robustez matemática trouxe flexibilidade estrutural do modelo quanto à sua abrangência espacial. A formulação matemática elaborada compreende e adapta-se facilmente a diferentes escalas de operação, podendo ser aplicada em malhas logísticas de nível estritamente municipal, estadual, nacional, ou em redes híbridas que integrem uma ou até três dessas esferas simultaneamente. Apesar das limitações listadas, a modelagem desenvolvida provou ser fortemente fidedigna ao problema logístico real.

Os resultados gerados pelo sistema foram extremamente precisos e informativos. O auxilio do Power BI permitiu uma interpretação visual clara dos dados gerados pelo algoritmo em Python. Por meio dos resultados, foi possível identificar rotas gargalo na rede logística, o percentual de uso dos modais, o total de doses movimentadas e o tempo exato ou aproximado usado para realizar toda a distribuição e entender o impacto da existência da intermediação estadual no processo.

Mesmo diante de um cenário de complexidade colossal, o Gurobi foi capaz de encontrar a solução ótima das instâncias de forma quase instantânea. Dessa forma o \textit{solver} revelou-se excepcionalmente poderoso.

A relevância deste estudo também se estende à sua notável adaptabilidade. A lógica de otimização para redistribuição em uma rede previamente conhecida pode ser reaproveitada para o transporte de outros insumos, tanto na área da saúde quanto em setores estratégicos, como o militar. A aplicação do modelo abrange desde a logística ágil de materiais de emergência até a gestão eficiente de materiais de fornecimento contínuo amparados por contratos previamente firmados.

Embora tenha sido feita simplificações, este trabalho comprova que uma formulação matemática coesa com a realidade e aliada a um algoritmo de ponta, prova ser uma ferramenta de gestão tática de alto nível e aplicável ao mundo real.

Este trabalho reafirma a importância da pesquisa operacional como recurso indispensável para a gestão da saúde pública. A transição de um planejamento de rotas empírico e manual para um sistema automatizado e guiado por dados cumpre o propósito maior de garantir a imunização da população de forma eficiente. O trabalho desenvolvido deixa um legado prático, escalável e cientificamente embasado para o enfrentamento de crises de pandemia, mitigação de desperdício de vacinas e preservação de vidas no mundo.
\bigskip